% Chapter 12 worked example. Included by ch12_resources.tex.
\section{Worked Example: Readmission-Risk Lifecycle Decision}
\label{sec:comp-ch12-worked-example}

\subsection{Purpose and Status}
This worked example demonstrates how the methods in Chapter~12 can be
translated into a reviewable institutional decision. All numerical inputs are
synthetic unless a source is identified explicitly. They are intended to show the
logic of the analysis, not to report empirical findings or establish a universal
threshold. Local use requires replacement of the assumptions with current,
versioned evidence and approval through the relevant clinical, managerial,
economic, legal, technical, and governance processes.

Following local validation and a controlled pilot, the hospital must decide whether to scale the readmission-risk intervention. The evidence dossier includes model performance, care-management outcomes, cost, subgroup effects, workload, vendor updates, and unresolved uncertainty. 

\subsection{Decision Problem and Comparators}
\textbf{Decision problem.} Should the intervention scale, continue with conditions, remain restricted, be redesigned, or be de-adopted?

The comparator set is deliberately broader than the existing pathway. This prevents
a new technology from receiving a lower evidentiary threshold than staffing,
workflow, shared-service, conventional digital, or no-change alternatives.

\begin{longtable}{@{}p{0.08\textwidth}p{0.84\textwidth}@{}}
\caption{Comparators for the Chapter 12 worked example}
\label{tab:comp-ch12-comparators}\\
\toprule
\textbf{Option} & \textbf{Comparator or strategy} \\
\midrule
\endfirsthead
\toprule
\textbf{Option} & \textbf{Comparator or strategy} \\
\midrule
\endhead
\bottomrule
\endlastfoot
1 & Continue limited pilot \\
2 & Scale with current threshold and capacity \\
3 & Scale with expanded care-management capacity \\
4 & Restrict to selected services or subgroups \\
5 & Return to the conventional discharge pathway \\
\end{longtable}

\subsection{Model and Decision Structure}
The analytical structure should follow the decision rather than the available
software. The team should map the causal pathway from intervention input to
information, human decision, action, outcome, resource consequence, and review.
The structure should also identify capacity constraints, uncertainty, subgroup
variation, implementation requirements, and events that would change the
recommendation.

A compact quantitative relationship used in this example is:
\begin{equation}
\text{Decision}=f(\text{evidence, value, feasibility, safeguards, ownership, reversibility})
\label{eq:comp-ch12-key-relationship}
\end{equation}
The equation is an organizing aid rather than a substitute for disaggregated evidence,
qualitative findings, or mandatory safeguards. Its variables and interpretation must
be defined in the local protocol.

\subsection{Synthetic Parameters and Evidence Sources}
Table~\ref{tab:comp-ch12-parameters} provides the base-case inputs. A
production analysis should add parameter distributions, correlations, structural
alternatives, data dates, system versions, and evidence-confidence ratings.

\begin{longtable}{@{}p{0.32\textwidth}p{0.22\textwidth}p{0.38\textwidth}@{}}
\caption{Synthetic parameters for the Chapter 12 worked example}
\label{tab:comp-ch12-parameters}\\
\toprule
\textbf{Parameter} & \textbf{Base value} & \textbf{Source or interpretation} \\
\midrule
\endfirsthead
\toprule
\textbf{Parameter} & \textbf{Base value} & \textbf{Source or interpretation} \\
\midrule
\endhead
\bottomrule
\endlastfoot
Pilot duration & 9 months & AI-HED record \\
Readmission relative risk & 0.91 & Local controlled evaluation \\
Incremental cost per discharge & EUR 94 & TCAIO record \\
Care-management capacity used & 93\% & Operations dashboard \\
Subgroup with calibration concern & Patients with incomplete social-risk data & EWAHE review \\
Material vendor update expected & Within 6 months & Change notice \\
\end{longtable}

\subsection{Step-by-Step Analysis}
The sequence below is intended to make the analysis reproducible and to ensure
that evidence, economics, governance, and implementation develop together.

\begin{longtable}{@{}p{0.07\textwidth}p{0.55\textwidth}p{0.30\textwidth}@{}}
\caption{Analytical sequence}
\label{tab:comp-ch12-analysis-steps}\\
\toprule
\textbf{Step} & \textbf{Required work} & \textbf{Decision record} \\
\midrule
\endfirsthead
\toprule
\textbf{Step} & \textbf{Required work} & \textbf{Decision record} \\
\midrule
\endhead
\bottomrule
\endlastfoot
1 & Use AI-HEVG to integrate problem relevance, data readiness, clinical validity, economic value, ethics, and organizational feasibility. & Evidence and responsible owner to be recorded \\
2 & Use AI-HED as the living decision record linking evidence, assumptions, decisions, and versions. & Evidence and responsible owner to be recorded \\
3 & Apply non-compensatory safeguards before considering aggregate scores. & Evidence and responsible owner to be recorded \\
4 & Assign accountable roles through HOAM and define monitoring, incidents, material change, renewal, and exit. & Evidence and responsible owner to be recorded \\
5 & Use GATED-AI, VALIDATE-AI, and other processes only for their defined operational roles. & Evidence and responsible owner to be recorded \\
6 & Reach a decision with conditions, monitoring thresholds, review dates, and stopping authority. & Evidence and responsible owner to be recorded \\
\end{longtable}

\subsection{Illustrative Outputs}
The outputs in Table~\ref{tab:comp-ch12-outputs} show how technical,
clinical, operational, economic, equity, and governance findings should be kept
visible rather than compressed into a single favourable or unfavourable score.

\begin{longtable}{@{}p{0.30\textwidth}p{0.24\textwidth}p{0.38\textwidth}@{}}
\caption{Illustrative model and decision outputs}
\label{tab:comp-ch12-outputs}\\
\toprule
\textbf{Output} & \textbf{Result} & \textbf{Interpretation} \\
\midrule
\endfirsthead
\toprule
\textbf{Output} & \textbf{Result} & \textbf{Interpretation} \\
\midrule
\endhead
\bottomrule
\endlastfoot
Clinical effect & Modest and uncertain & Direction favourable; magnitude below initial expectation \\
Economic value & Conditional & Dependent on capacity expansion and avoided admissions \\
Equity concern & Material but potentially remediable & Incomplete social-risk data affect prioritisation \\
Lifecycle concern & Upcoming material update & Requires version-specific reassessment \\
Preferred decision & Continue with conditions; do not scale fully & Evidence and capacity not yet sufficient \\
\end{longtable}

\subsection{Sensitivity, Uncertainty, and Subgroups}
At minimum, the completed analysis should vary the parameters most likely to
change the decision, test at least one credible alternative model structure, and
report subgroup results separately where performance, access, response, burden,
or opportunity cost may differ. Deterministic analysis should identify thresholds;
probabilistic analysis should be used where credible parameter distributions and
correlations are available. Scenario analysis is preferable when structural or
future uncertainty cannot be represented defensibly by one probability model.

The record should distinguish uncertainty that can be reduced through evidence,
uncertainty that can be managed through safeguards, and uncertainty that remains
too consequential to accept. Missing evidence should not be represented as an
average or neutral value without justification.

\subsection{Interpretation and Recommendation}
Continue the intervention within its current scope while addressing social-risk data, care-management capacity, and the forthcoming material update. Full scale should require renewed validation and a revised economic and equity assessment.

The recommendation is conditional rather than permanent. It applies only to the
specified population, setting, intervention configuration, evidence cut-off, and
version. The following conditions should be incorporated into the approval,
contract, implementation, monitoring, or renewal record:

\begin{itemize}
 \item AI-HED is updated with the exact deployed version, evidence cut-off, and decision rationale.
 \item monitoring includes workload, outcomes, subgroup access, cost, incidents, and trust.
 \item the material update is assessed before deployment and rollback remains available.
 \item the accountable authority can restrict or retire use without vendor approval.
 \item renewal is based on current evidence and realized value rather than prior investment.
\end{itemize}

\subsection{Decision Statement}
\begin{quote}
For the specified decision problem and comparator set, the institution should
\emph{[proceed / proceed with conditions / redesign / pilot / restrict / defer /
replace / retire]}. The decision applies to \emph{[population, setting, version,
and evidence period]}, is owned by \emph{[accountable authority]}, and will be
reviewed on \emph{[date]} or earlier if \emph{[trigger events]} occur. The
evidence, assumptions, unresolved uncertainty, mandatory safeguards, monitoring
thresholds, and exit arrangements are recorded in the accompanying dossier.
\end{quote}
